\subsection{Short connectors for every terminal pair \((2,\allowbreak b)\)}\label{proof:Q5-C085}
On the original nonempty finite probability space, let \(K\) be real and antisymmetric with \(|K|\le\allowbreak 1\), and put
\[
 Jv=(Ev)1,\quad A=J+K,\quad T=J-K=A^*,\quad
 g=K1,\quad p=1+g,\quad q=1-g,
\]
\[
 f_n=A^n1,\quad \bar f_n=T^n1,\quad h_n=Tf_{n-1},\qquad
 \delta_j=\|K^j1\|_2^2,\quad d=\delta_1,
\]
\[
 \theta=13/32,\quad\kappa=19/32,\quad
 r=Kg,\quad s=Kr,\quad t=Ks.
\]
All adjoints, norms, spectral measures and expectations use the same original \(\mu\).
Choose the terminal directions by the original integer lengths:
\[
\begin{array}{c|c}
\text{terminal and connector}&F_{b,L}\\ \hline
b\ge3,\ L=1&h_2f_b\\
b=3,\ L=2&h_2f_3\\
b\ge4,\ L=2&f_2h_b\\
b\ge3,\ L\ge3&h_2f_b
\end{array}
\]
Each \(f_n\) arm points entirely away from the terminal.
Each \(h_n\) arm has its first edge toward the terminal and its remaining edges away from its first neighbour toward the leaf.
All \(L\) connector edges point from the terminal toward the receiving root, whose original unit arm points outward.
The actual message is \(H=\allowbreak pT^LF_{b,\allowbreak L}\); these directions are fixed before the host is given.

The short-connector bounds are
\[
 E[pT^2(h_2f_3)]^8\le e^{-d/25},                         \tag{Q85.0.1}\label{eq:Q85:0.1}
\]
\[
 E[pT^2(f_2h_4)]^8\le e^{-d/3},\qquad
 E[pT^3(h_2f_4)]^8\le e^{-d/3},                        \tag{Q85.0.2}\label{eq:Q85:0.2}
\]
\[
 E[pT^2(f_2h_b)]^8\le e^{-d/5},\qquad
 E[pT^3(h_2f_b)]^8\le e^{-d/5}\quad(b\ge5).             \tag{Q85.0.3}\label{eq:Q85:0.3}
\]
The high-domain part of the second bound in \eqref{eq:Q85:0.3} uses the terminal-mean theorem in Section~\ref{proof:Q5-C086}; its use is stated below.
Together with the unit and long-connector rows, these estimates give
\[
 E H^8\le Q(d):=1-d/20000.                              \tag{Q85.0.4}\label{eq:Q85:0.4}
\]
The ordinary row \(F_{b,\allowbreak 2}=\allowbreak f_2h_b\) for \(b\ge\allowbreak 4\) has its stated direction.
The \(f_2\)-reference fourth bound for \(h_2f_b\) is a separate input on that actual branch.

\paragraph{Q85.1 Finite framework and local inputs}
The finite calculus in Section~\ref{proof:Q5-C081} gives
\[
 0\le d\le1/3,\quad\|K\|^2\le\theta,\quad
 d^2\le \delta_2\le\theta d,\quad \delta_{j+1}\le\theta \delta_j,
\]
\[
 |Kv(x)|\le\|v\|_2,\quad
 EK^{2j}1=(-1)^j \delta_j,\quad EK^{2j+1}1=0.               \tag{Q85.1.1}\label{eq:Q85:1.1}
\]
After deleting zero-mass atoms, \(p(x),\allowbreak q(x)\ge\allowbreak \mu(x)>\allowbreak 0\).
Actual tree messages are positive: conditional on the root label, assigning every other label that same value has positive mass and makes every edge factor one.
At \(d=\allowbreak 0\), \(g=\allowbreak 0\) and successive leaf integration gives every actual message one, so all moment bounds hold directly.
Subsequent divisions by \(d\) use \(d>\allowbreak 0\).

The finite positive convolution is
\[
 S_n:=\sum_{j=0}^nK^j1
 =f_n+\sum_{j=3}^n\beta_jf_{n-j},\quad
 \beta_1=\beta_2=0,
\]
\[
 \beta_j=\sum_{i=1}^{\lfloor(j-1)/2\rfloor}(-1)^{i-1}\delta_i,
 \qquad \tau:=d-\delta_2\le\beta_j\le d\quad(j\ge3).        \tag{Q85.1.2}\label{eq:Q85:1.2}
\]
Its complete reversal uses \(K'=\allowbreak -K,\allowbreak g'=\allowbreak -g\).
All finite path messages in discarded convolution terms are nonnegative.

We use the signed coefficient-\(17\) exponential estimate in Section~\ref{proof:Q5-C081}, the retained-envelope formulas and full-frequency bounds of Section~\ref{proof:Q5-C077}, and the corresponding \(b=\allowbreak 4\) parameter bounds of Section~\ref{proof:Q5-C073}.
The arithmetic-mean input for \(h_4+f_v\), for every finite \(v\ge\allowbreak 5\), is the one in Section~\ref{proof:Q5-C067}; its needed formula is restated in \eqref{eq:Q85:6.4}'s derivation.
The ordinary unit row is \( \E[pT(h_2f_b)]^8\le\allowbreak 1-d/500\) for every \(b\ge\allowbreak 3\).
The remaining long rows and the high-domain actual terminal mean are supplied by Section~\ref{proof:Q5-C086}.

\paragraph{Q85.2 A variance-sensitive exponential comparison}\label{calc:Q85:variance-exp}
Let \(W\) be finite and real with \(|W|\le\allowbreak  U\) and
\(\E W^2\le\allowbreak \theta U^2\). For every \(s\ge\allowbreak 0\), retaining the actual mean of either sign,
\[
 E e^{sW}\le
 \exp\left\{sEW+\frac{s^2}{8}(U^2+3EW^2)\right\}.    \tag{Q85.2.1}\label{eq:Q85:2.1}
\]
In particular, for every real original vector \(V\),
\[
 E e^{8KV}\le
 \exp\{8E(KV)+8\|V\|^2+24\|KV\|^2\}.                \tag{Q85.2.2}\label{eq:Q85:2.2}
\]
The cases \(U=0\) and constant-variable endpoints are immediate.
Otherwise set \(h=\E W^2/U^2\), \(m=\E W\),
\(v=\E(W-m)^2\), and \(M=U-m\).
Apply \eqref{eq:oc:hermite-matching} to \(f(w)=e^{sw}\):
\[
 E e^{sW}\le e^{sm}E e^{sL(B-\pi)},\qquad
 L=M+v/M,\quad \pi=\frac{v}{M^2+v},                   \tag{Q85.2.3}\label{eq:Q85:2.3}
\]
where \(B\) is Bernoulli with mean \(\pi\).

When \(m\ge\allowbreak 0\), write \(w=\allowbreak m/U\). Then
\[
 L/U=\frac{1+h-2w}{1-w}\le1+h.
\]
The centered Bernoulli logarithmic moment has tilted second derivative at most \(L^2/4\).
Integrating twice gives coefficient \(L^2/8\), at most
\((1+3h)U^2/8\), since \((1+h)^2\le\allowbreak 1+3h\) for \(0\le\allowbreak  h\le\allowbreak 1\).

When \(m\le\allowbreak 0\), put \(u=\allowbreak -m/U,\allowbreak t=\allowbreak 1+u\) and
\[
 a=\frac{h-u^2}{(1+u)^2}\in[0,h],\qquad
 L/U=t(1+a),\quad y=1-2\pi=\frac{1-a}{1+a}.          \tag{Q85.2.4}\label{eq:Q85:2.4}
\]
The centered Bernoulli estimate is
\[
 E e^{z(B-\pi)}\le
 \exp\left\{\frac{1-2\pi}{4\log((1-\pi)/\pi)}z^2\right\}.
                                                               \tag{Q85.2.5}\label{eq:Q85:2.5}
\]
This is \eqref{eq:Q102:3.5}, including its continuous endpoints.

Use \(\log((1+y)/(1-y))\ge\allowbreak 2(y+y^3/3)\).
It suffices to bound \(L^2/[8(1+y^2/3)]\).
After division by \(t^2\), the desired inequality becomes
\[
 \frac{3(1+a)^4}{4(1+a+a^2)}
 \le \frac{1+3u^2}{(1+u)^2}+3a.                     \tag{Q85.2.6}\label{eq:Q85:2.6}
\]
Subtracting \(3a\) on the left yields
\[
 \frac{3(1+a^2)^2}{4(1+a+a^2)}\le\frac34,
\]
because \(a^3+a\le\allowbreak \theta^3+\theta<\allowbreak 1\). On the right,
\[
 \frac{1+3u^2}{(1+u)^2}-\frac34
 =\frac{(3u-1)^2}{4(1+u)^2}\ge0.
\]
This proves \eqref{eq:Q85:2.1}.
The original row and spectral bounds, with \(U=\allowbreak \|V\|\) and \(W=\allowbreak KV\), prove \eqref{eq:Q85:2.2}.
Degenerate Bernoulli endpoints follow directly or by continuity.
The coefficient-\(17\) estimate and \eqref{eq:Q85:2.2} can both be applied to the same primitive.

\paragraph{Q85.3 The exceptional row \(b=\allowbreak 3,\allowbreak L=\allowbreak 2\)}
Use \(F=\allowbreak h_2f_3\). Its arithmetic mean is
\[
 h_2+f_3=2-d+s=2m+s,\quad m=1-d/2,\quad k=1/m\in[1,6/5].
\]
Put \(K'=\allowbreak -K,\allowbreak g'=\allowbreak K'\one,\allowbreak a=\allowbreak k\delta_2,\allowbreak \delta=\allowbreak k^2\delta_3/2\) and
\(\eta=\allowbreak a-\delta\).
Positive two-step propagation of \(F\le\allowbreak (m+s/2)^2\) gives
\[
 0<X=T^2F\le m^2 Z,
 \quad Z=1-\eta+(1+\delta+K'-kK'^4)g'.               \tag{Q85.3.1}\label{eq:Q85:3.1}
\]
This is the two-step envelope obtained by \(T(s^2)\le\allowbreak 2\delta_3\).
Set
\[
 V=\delta1+g'-kK'^3g',\qquad N=K'V.
\]
Then \(p+Z=\allowbreak 2-\eta+N\), with \(p,\allowbreak Z\ge\allowbreak 0\), so
\[
 (pX)^8\le m^{16}[1+(N-\eta)/2]^{16}
 \le m^{16}\exp(8N-8\eta).                          \tag{Q85.3.2}\label{eq:Q85:3.2}
\]
The base is the actual arithmetic mean.
Its zero endpoint is included by
\((1+x)^{16}\le\allowbreak  e^{16x}\) for \(x\ge\allowbreak -1\).
Skew orthogonality gives
\[
 EN=-d,
\]
\[
 \|V\|^2=d+k^2\delta_4-2\delta k \delta_2+\delta^2
 \le d+k^2\delta_4,
\]
\[
 \|K'V\|^2=\delta_2+k^2\delta_5-2\delta k \delta_3+\delta^2d
 \le \delta_2+k^2\delta_5.                                    \tag{Q85.3.3}\label{eq:Q85:3.3}
\]
The discarded terms are nonpositive because
\(\delta\le\allowbreak 2k\delta_2\) and \(\delta d\le\allowbreak 2k\delta_3\).
Both follow from \(\delta=\allowbreak k^2\delta_3/2,\allowbreak k\le\allowbreak 6/5,\allowbreak d\le\allowbreak 1/3\) and
\(\delta_3\le\allowbreak \theta\delta_2\).

Using the coefficient-\(17\) estimate in \eqref{eq:Q85:3.2}, and
\(16\log m\le\allowbreak -8d\), yields
\[
 \log E(pX)^8
 \le d-8k \delta_2+4k^2\delta_3+17k^2\delta_4
 \le d-(78/25)\delta_2.                                  \tag{Q85.3.4}\label{eq:Q85:3.4}
\]
Indeed \(4\theta+17\theta^2=\allowbreak 4537/1024<\allowbreak 9/2\), and
\(-8k+(9/2)k^2\) increases on \([1,\allowbreak 6/5]\), with maximum \(-78/25\).
Using \eqref{eq:Q85:2.2} instead gives
\[
 \log E(pX)^8\le
 -8d+(24-8k)\delta_2+4k^2\delta_3+8k^2\delta_4+24k^2\delta_5
 \le-8d+(108/5)\delta_2.                                 \tag{Q85.3.5}\label{eq:Q85:3.5}
\]
Here \(4\theta+8\theta^2+24\theta^3<\allowbreak 5\): replacing \(\theta\) by \(5/12\) already gives \(115/24<\allowbreak 5\).
The quadratic \(24-8k+5k^2\) increases on \([1,\allowbreak 6/5]\) and has maximum \(108/5\).
If \(\delta_2/d\ge\allowbreak 1/3\), \eqref{eq:Q85:3.4} is at most \(-d/25\).
If \(\delta_2/d\le\allowbreak 1/3\), \eqref{eq:Q85:3.5} is at most \(-4d/5\).
Both include equality, proving \eqref{eq:Q85:0.1} on the full domain.

\paragraph{Q85.4 Two actual signs of the retained square envelope}
For \(b=\allowbreak 4\), let \(m=\allowbreak 1-d,\allowbreak R=\allowbreak s+t-dg\). The finite formulas give
\[
 h_2+f_4=2m+R,\qquad f_2+h_4=2m-R.                  \tag{Q85.4.1}\label{eq:Q85:4.1}
\]
For every \(b\ge\allowbreak 5\) on \(0<\allowbreak d\le\allowbreak 1/5\), use the retained coefficients with \(m=\allowbreak 1-9d/8\).
For \(b=\allowbreak 5\), take
\[
 c_1=d,\quad c_2=d/4+\delta_2,\quad c_3=0.
\]
For \(b\ge\allowbreak 6\), set
\[
 \lambda=
 \frac{9d/4-\beta_b-\beta_{b-1}}
 {\beta_{b-2}+(1-d)\beta_{b-3}},\quad
 c_1=\beta_{b-1},\ c_2=\lambda\beta_{b-2},\ c_3=\lambda\beta_{b-3}.
                                                               \tag{Q85.4.2}\label{eq:Q85:4.2}
\]
The denominator is positive and \(0\le\allowbreak \lambda\le\allowbreak 1\):
the full retained constant is at least
\((4-d)\tau\ge\allowbreak 361d/160>\allowbreak 9d/4\).
Thus every \(c_i\) is bounded by its actual positive-convolution coefficient.
Put
\[
 C=c_1+c_2+c_3,\quad D=c_2+c_3,\quad E_*=c_3,
\]
\[
 R=\sum_{j=3}^b K^j1-Cg-Dr-E_*s.
\]
In the positive convolution for \(f_b\), retain
\(\beta_b,\allowbreak c_1f_1,\allowbreak c_2f_2,\allowbreak c_3f_3\). This gives \(h_2+f_b\le\allowbreak 2m+R\).
For the opposite pair, apply \(T\) to the positive convolution for \(f_{b-1}\), retaining \(c_1q,\allowbreak c_2h_2,\allowbreak c_3h_3\). Since
\[
 h_3=(1-d)-g-r-s,
\]
the constant loss is again
\(\beta_b+c_1+c_2+(1-d)c_3=\allowbreak 9d/4\), and the remaining polynomial is \(-R\).
Thus
\[
 0<h_2+f_b\le2m+R,\qquad
 0<f_2+h_b\le2m-R.                                  \tag{Q85.4.3}\label{eq:Q85:4.3}
\]
The actual products are bounded by \((m+R/2)^2\) and \((m-R/2)^2\), respectively.

Write \(K'=\allowbreak -K,\allowbreak g'=\allowbreak -g\).
On its rotation blocks \(z=\allowbreak i\sqrt x\),
\(R=\allowbreak (\mathsf A(x)-\mathsf B(x)z)g'\).
All unlabelled integrals below use the finite positive spectral measure of
\(-K'^2\) at \(g'\), with total mass \(d\).
For \(b=\allowbreak 5\),
\[
 \mathsf A=d+D+x-x^2,\qquad\mathsf B=D+x,
 \quad D=d/4+\delta_2.                                   \tag{Q85.4.4}\label{eq:Q85:4.4}
\]
For \(b\ge\allowbreak 6\), let \(N=\allowbreak b-2\ge\allowbreak 4\) and \(G_N(z)=\allowbreak a_N+zb_N\). Then
\[
 \mathsf A=C+x(a_N-E_*),\quad
 \mathsf B=D+x b_N,                                  \tag{Q85.4.5}\label{eq:Q85:4.5}
\]
\[
 1-x\le a_N,b_N\le1-x+x^2,\quad a_N\ge(1-x)b_N.
\]
These are finite alternating geometric sums. The coefficients and the global
\(\delta_2\) in \(c_i\) stay fixed while \(x\) varies.
Put \(k=\allowbreak 1/m\) and
\[
 V_R=ER^2=\int(\mathsf A^2+x\mathsf B^2),\quad
 D_0=ER=\int\mathsf B,\quad a_1=EKR=\int\mathsf A,
\]
\[
 \rho=-E(rR)=\int x\mathsf B,\quad
 U_R=\|KR\|^2=\int x(\mathsf A^2+x\mathsf B^2),
\]
\[
 \delta=kD_0+k^2V_R/2,\quad a=ka_1,\quad
 \varepsilon=\delta-a.                              \tag{Q85.4.6}\label{eq:Q85:4.6}
\]
The retained-envelope estimates give
\[
 \mathsf A,\mathsf B\ge0,\quad
 \mathsf A-(1-x)\mathsf B\ge c_1+c_2x,\quad c_1\ge\tau,
\]
\[
 \sqrt{\mathsf A^2+x\mathsf B^2}\le u_0=3/8+9d/5,
 \quad \mathsf B\le z_0=\theta+13d/16,                \tag{Q85.4.7}\label{eq:Q85:4.7}
\]
\[
 |\varepsilon|\le3d/5,\quad \delta\le6d/5,\quad ku_0\le1,
 \quad a_1-D_0\le d^2+\theta^3d.                     \tag{Q85.4.8}\label{eq:Q85:4.8}
\]
For the last inequality, \eqref{eq:Q85:4.5} gives
\(\mathsf A-\mathsf B=\allowbreak c_1-E_*x+x(a_N-b_N)\);
the first possible positive finite end term is \(x^3\).
For \(b=\allowbreak 5\), the difference integrates to \(d^2-\delta_3\).
Multiplying the same finite observation by \(x\) also gives
\[
 D_3-\rho:=\int x(\mathsf A-\mathsf B)
 \le d \delta_2+\delta_5.                                     \tag{Q85.4.9}\label{eq:Q85:4.9}
\]
For \(b=\allowbreak 4\) on \(0<\allowbreak d\le\allowbreak 1/4\), the corresponding exact quantities are
\[
 D_0=\delta_2,\quad a_1=d^2+\delta_2,\quad \rho=\delta_3,
\]
\[
 V_R=d^3+2d\delta_2+\delta_3+\delta_4,\quad
 U_R=d^2\delta_2+2d\delta_3+\delta_4+\delta_5,
\]
\[
 \delta=k\delta_2+k^2V_R/2,\quad
 \varepsilon=-kd^2+k^2V_R/2,
 \quad |\varepsilon|\le d/6,\quad\delta\le25d/24.     \tag{Q85.4.10}\label{eq:Q85:4.10}
\]
The last two estimates apply to this same \(R,\allowbreak m\).

\paragraph{Q85.5 Full-frequency contraction of the retained envelope}\label{calc:Q85:retained-contraction}
For every \(b\ge\allowbreak 5,\allowbreak 0<\allowbreak d\le\allowbreak 1/5\) and every original frequency
\(0\le\allowbreak  x\le\allowbreak \theta\),
\[
 \boxed{\mathsf A(x)^2+x\mathsf B(x)^2\le2m\mathsf B(x).} \tag{Q85.5.1}\label{eq:Q85:5.1}
\]
Consequently,
\[
 V_R\le2mD_0,\qquad kU_R\le2\rho.                  \tag{Q85.5.2}\label{eq:Q85:5.2}
\]
First suppose \(b\ge\allowbreak 6\), and let \(f=\allowbreak x-x^2+x^3,\allowbreak \ell=\allowbreak x-x^2\).
Equations \eqref{eq:Q85:4.2}--\eqref{eq:Q85:4.5} give
\[
 \mathsf A\le D+d+f,\quad D+\ell\le\mathsf B\le D+f,
\]
\[
 d/4\le D\le(17/16)d+d^2.                           \tag{Q85.5.3}\label{eq:Q85:5.3}
\]
Indeed
\(D=\allowbreak 9d/4-\beta_b-\beta_{b-1}+dE_*\), and these two global bounds follow from
\(\tau\le\allowbreak \beta_j\le\allowbreak  d,\allowbreak E_*\le\allowbreak  d\). It is enough to prove
\[
 Q=(1+x)(D+f)^2+2d(D+f)+d^2-2m(D+\ell)\le0.         \tag{Q85.5.4}\label{eq:Q85:5.4}
\]
For fixed \(d\), \(Q\) is convex in \(D\). It is also convex in \(x\).
To see the latter, use
\[
 (1+x)f=x+x^4,
\]
\[
 (1+x)f^2=x^2-x^3+x^4+x^5-x^6+x^7=:P(x).
\]
Here \(P''\ge\allowbreak 5/4\), since
\(2-6x+12x^2\ge\allowbreak 5/4\), \(20x^3-30x^4\ge\allowbreak 0\), and the remaining terms are nonnegative.
The potentially negative term \(2df''\ge\allowbreak -4d\) is paid by the contribution \(4m\).
Thus \(Q_{xx}>\allowbreak 0\) on the whole rectangle.
It suffices to use \(x=\allowbreak 0,\allowbreak \theta\) and the two \(D\) endpoints while retaining the full interval of \(d\).

At \(x=\allowbreak 0,\allowbreak D=\allowbreak d/4\),
\[
 Q=d[(17/8)d-1/2]\le-3d/40.
\]
At \(x=\allowbreak 0,\allowbreak D=\allowbreak (17/16)d+d^2\),
\[
 Q/d=-17/8+(1189/256)d+(51/8)d^2+d^3<0.
\]
For the last inequality, increase the three positive coefficients to \(5,\allowbreak 7,\allowbreak 1\) and use \(d\le\allowbreak 1/5\).
At \(x=\allowbreak \theta\), use
\(f(\theta)<\allowbreak 5/16\), \(6/25\le\allowbreak \ell(\theta)\le\allowbreak 1/4\) and \(1+\theta<\allowbreak 3/2\).
For \(D=\allowbreak d/4\),
\[
 Q\le75/512-12/25+(59/64)d+(69/32)d^2
 \le-161/2560.                                      \tag{Q85.5.5}\label{eq:Q85:5.5}
\]
For \(D=\allowbreak (17/16)d+d^2\), the linear coefficient is at most \(-15/4096\).
The quadratic, cubic and quartic coefficients are less than \(6,\allowbreak 15/2,\allowbreak 3/2\), respectively, so
\[
 Q\le75/512-12/25+6d^2+(15/2)d^3+(3/2)d^4
 \le-9957/320000.                                    \tag{Q85.5.6}\label{eq:Q85:5.6}
\]
Before enlargement the quadratic coefficient is \(49005/8192<\allowbreak 6\), and the cubic coefficient is \(1853/256<\allowbreak 15/2\).
The linear comparison is
\((1277/256)(5/16)-17/8+9/16=\allowbreak -15/4096\).
Every remaining positive polynomial was bounded at \(d=\allowbreak 1/5\).
Convexity proves the intermediate frequencies and \(D\) values.

For \(b=\allowbreak 5\),
\(\mathsf A=\allowbreak D+d+x-x^2,\allowbreak \mathsf B=\allowbreak D+x\), with
\(d/4\le\allowbreak  D\le\allowbreak (21/32)d\).
Because \(\mathsf A\ge\allowbreak 0\), discard its \(-x^2\) inside the positive square.
It suffices to prove
\[
 Q=(1+x)(D+x)^2+2d(D+x)+d^2-2m(D+x)\le0.             \tag{Q85.5.7}\label{eq:Q85:5.7}
\]
This \(Q\) is convex in \(D\) and \(x\); its second \(x\)-derivative is
\(2(1+2D)+6x>\allowbreak 0\).
At \(x=\allowbreak 0\), the lower \(D\) endpoint was proved above, and the upper gives
\(Q/d=\allowbreak -21/16+(4321/1024)d<\allowbreak 0\).
At \(x=\allowbreak \theta\), write \(D=\allowbreak t_0d\), \(t_0=\allowbreak 1/4\) or \(21/32\).
The constant is \(-19019/32768\le\allowbreak -9/16\), the linear coefficient is at most \(13/8\), and the quadratic coefficient is at most \(9/2\).
For the quadratic bound, use \(1+\theta<\allowbreak 3/2,\allowbreak t_0<\allowbreak 2/3\);
for the linear bound, use \(2\theta+2\theta^2-2<\allowbreak -3/4,\allowbreak t_0\ge\allowbreak 1/4\).
Therefore
\[
 Q\le-9/16+(13/8)d+(9/2)d^2\le-23/400.
\]
This proves \eqref{eq:Q85:5.1} for every finite \(b\ge\allowbreak 5\).
Integrate against the original spectral measure, first directly and then after multiplication by \(x\), to obtain \eqref{eq:Q85:5.2}.
The global \(\delta_2\) is constant in this frequency calculation.

\paragraph{Q85.6 Every \(b\ge\allowbreak 5,\allowbreak L=\allowbreak 2\), in the reversed-pair direction}
For \(0<\allowbreak d\le\allowbreak 1/5\), take \(F=\allowbreak f_2h_b\) and use its actual square
\((m-R/2)^2\) from \eqref{eq:Q85:4.3}.
In the two-step Duhamel calculation write
\[
 \delta_-=-kD_0+k^2V_R/2,\quad a_-=-ka_1,
 \quad\varepsilon_-=\delta_--a_-.
\]
Positive propagation gives
\[
 0<T^2F/m^2\le Z=\bar f_2-kK'^2R+ka_1+\delta_-q.
\]
This follows by expanding that square and using \(T(R^2)\le\allowbreak 2V_R\).
In particular its right side is nonnegative. It satisfies
\[
 p+Z=2+\varepsilon_-+N,
 \quad N=K'V_-,\quad V_-=\delta_-1+g'-kK'R.
\]
The exact moments are
\[
 EN=-d+k\rho,
\]
\[
 \|V_-\|^2=d-2k\rho+k^2U_R+\delta_-^2-2\delta_-a_-.
                                                               \tag{Q85.6.1}\label{eq:Q85:6.1}
\]
Equation \eqref{eq:Q85:5.2} gives \(\delta_-\le\allowbreak 0\) and \(kU_R\le\allowbreak 2\rho\).
Also \(a_1\ge\allowbreak \kappa D_0+\tau d\) implies \(D_0\le\allowbreak 2a_1\), so
\(\delta_-\ge\allowbreak 2a_-\).
Thus the entire constant-coordinate term in \eqref{eq:Q85:6.1} is nonpositive and
\[
 \|V_-\|^2\le d.
\]
AM--GM and the coefficient-\(17\) estimate yield
\[
 \log E[pT^2F]^8
 \le16\log m+9d+8(\varepsilon_-+k\rho).              \tag{Q85.6.2}\label{eq:Q85:6.2}
\]
For the same \(R\),
\[
 \varepsilon_-+k\rho
 \le \{k(d+\theta^3)+(k^2/2)u_0^2+k\theta z_0\}d.
\]
All three terms increase with \(d\) on \([0,\allowbreak 1/5]\), and their endpoint comparisons are
\[
 43753/126976<7/20,\quad
 21609/48050<9/20,\quad
 1183/3968<3/10.
\]
Their sum is therefore less than \(11/10\).
Since \(16\log m\le\allowbreak -18d\), \eqref{eq:Q85:6.2} is at most \(-d/5\).
This proves the low-domain row for every finite \(b\ge\allowbreak 5\).

For \(1/5\le\allowbreak  d\le\allowbreak 1/3\), original adjunction gives
\[
 EF=E(f_2h_b)=E(f_3f_{b-1})=E(h_4f_{b-2}).            \tag{Q85.6.3}\label{eq:Q85:6.3}
\]
For \(b=\allowbreak 5\), this is
\(\E[f_3Af_3]=\allowbreak \E[f_3Tf_3]=\allowbreak (\E f_3)^2=\allowbreak (1-2d)^2\).
For \(b=\allowbreak 6\), it is
\(\E(h_4f_4)\le\allowbreak \mu_3^2=\allowbreak (1-2d)^2\), since \(h_4+f_4=\allowbreak 2\mu_3\) pointwise.
For \(b\ge\allowbreak 7\), set \(v=\allowbreak b-2\ge\allowbreak 5\). The arithmetic-mean input gives
\[
 h_4+f_v\le2-13d/4+R_{4v},\quad
 R_{4v}=\sum_{j=5}^vK^j1.
\]
The finite alternating means start with a zero odd term and a negative even term, so
\(\E R_{4v}\le\allowbreak 0\). Its norm satisfies
\[
 ER_{4v}^2\le(1+\theta)\delta_5\le(1+\theta)\theta^4d.
\]
The finite \(G_N\) norm is at most one except at \(N=\allowbreak 2\), where its square is at most
\(1+\theta\); \(N=\allowbreak 1\) is included.
Since \(1-13d/8>\allowbreak 0\),
\[
 E(h_4f_v)\le(1-13d/8)^2+\frac{(1+\theta)\theta^4}{4}d
 \le(1-13d/8)^2+d/80.                                \tag{Q85.6.4}\label{eq:Q85:6.4}
\]
The coefficient comparison already follows from \(\theta<\allowbreak 5/12\).
The last quadratic decreases on \([1/5,\allowbreak 1/3]\), with maximum
\(733/1600<\allowbreak 11/24\). The \(b=\allowbreak 5,\allowbreak 6\) maximum \(9/25\) is smaller.
Thus \(TF\le\allowbreak 2\E F\le\allowbreak 11/12\) at every atom.
Positivity gives \(T^2F\le\allowbreak (11/12)q\), whence
\[
 E[pT^2F]^8\le(11/12)^8\le e^{-2/3}\le e^{-d/5}.      \tag{Q85.6.5}\label{eq:Q85:6.5}
\]
The same direction applies on both domains, proving the first inequality in \eqref{eq:Q85:0.3}.

\paragraph{Q85.7 The two short \(b=\allowbreak 4\) rows}
For \(L=\allowbreak 2,\allowbreak F=\allowbreak f_2h_4,\allowbreak 0<\allowbreak d\le\allowbreak 1/4\), use \eqref{eq:Q85:4.10}.
In the reversed square,
\(\delta_-=\allowbreak -k\delta_2+k^2V_R/2\le\allowbreak 0\), because
\[
 V_R/\delta_2\le3d+\theta+\theta^2\le2(1-d).
\]
Also \(\varepsilon_-=\allowbreak kd^2+k^2V_R/2\ge\allowbreak 0\), so
\(\delta_-\ge\allowbreak  a_-\) and the constant norm contribution is nonpositive.
The two-step calculation in \eqref{eq:Q85:6.1} gives
\[
\begin{aligned}
 \log EH^8\le{}&16\log m+9d+8kd^2\\
 &+k^2[4d^3+(8d+17d^2)\delta_2+(4+34d)\delta_3+21\delta_4+17\delta_5]
 -26k\delta_3.                                           
\end{aligned}
\tag{Q85.7.1}\label{eq:Q85:7.1}
\]
Group \(\delta_3,\allowbreak \delta_4,\allowbreak \delta_5\) before estimating their signs.
Since \(4+21\theta+17\theta^2<\allowbreak 31/2\), the total coefficient of \(\delta_3\) is at most
\(k^2(31/2+34d)-26k\le\allowbreak 24k^2-26k\le\allowbreak 8\).
Bound the remaining positive terms at \(d=\allowbreak 1/4,\allowbreak k=\allowbreak 4/3\).
Using \(16\log m\le\allowbreak -16d\), the result divided by \(d\) is
\[
 -7+8/3+4/9+49\theta/9+8\theta^2=-137/384<-1/3.       \tag{Q85.7.2}\label{eq:Q85:7.2}
\]

For \(L=\allowbreak 3,\allowbreak F=\allowbreak h_2f_4\), use the canonical parameters in
\eqref{eq:Q85:4.10}. The three-step arithmetic-mean primitive is
\[
 V_3=\varepsilon1+(1+\delta)g'+K'g'+kK'^2R.
\]
Its nonconstant spectral coordinates are
\[
 1+\delta-kdx-kx^2,\qquad 1+kx^2.
\]
The scalar constant plus the signed mean of \(K'V_3\) is
\(-2d-kd^2+kV_R/2\). Thus
\[
 \log EH^8\le16\log m-16d-8kd^2+4kV_R+17\|V_3\|^2. \tag{Q85.7.3}\label{eq:Q85:7.3}
\]
In the nonconstant norm, group \(\delta_3,\allowbreak \delta_4,\allowbreak \delta_5,\allowbreak \delta_6\).
Their coefficient is at most
\[
 k[-2\kappa+k((d+\theta)^2+\theta^3)]<0,
\]
because
\((d+\theta)^2+\theta^3\le\allowbreak 16309/32768<\allowbreak 1/2\) and \(k\le\allowbreak 4/3\).
The constant-vector contribution is
\(\varepsilon^2-2\varepsilon(d+k\delta_3)
\le(1001/2304)d^2<d^2/2\).
The parameter bounds on \(\delta,\allowbreak \varepsilon\) give
\[
 \|V_3\|^2/d\le1+\theta+(85/48)d+(625/576)d^2.        \tag{Q85.7.4}\label{eq:Q85:7.4}
\]
The coefficient \(1-2k(1+\delta)d\) of \(\delta_2\) is nonnegative:
its minimum from those parameter bounds is \(23/144\).
Thus \(\delta_2\le\allowbreak \theta d\) can be used in this term.
Use \(16\log m\le\allowbreak -16d-8d^2,\allowbreak k\ge\allowbreak 1\) and
\(4kV_R/d\le\allowbreak 16309/6144\), and bound the remaining positive powers at \(d=\allowbreak 1/4\).
Equation \eqref{eq:Q85:7.3}, divided by \(d\), is at most
\[
 -259/32+677/192+10625/9216+16309/6144
 =-14015/18432<-3/4.                                 \tag{Q85.7.5}\label{eq:Q85:7.5}
\]

For the actual bare path, Corollary~\ref{cor:unit-square:two-step-path}
gives
\[
 \E(p\bar f_2)^8\le82/15\qquad(1/4\le d\le1/3).
 \tag{Q85.7.6}\label{eq:Q85:7.6}
\]
On the adjacent interval it gives
\[
 \E(p\bar f_2)^8\le35/8\qquad(1/5\le d\le1/4).
 \tag{Q85.7.7}\label{eq:Q85:7.7}
\]
Both statements use the same actual \(r=K^2\one\) and its compatible
mean and second moment, with the proof now shared at the unit root.

For the \(b=\allowbreak 4\) reversed pair on \(d\ge\allowbreak 1/4\),
\[
 E(f_2h_4)=Ef_3^2=1-3d+3d^2-\delta_2+\delta_3
 \le1-3d+(77/32)d^2\le205/512                       \tag{Q85.7.8}\label{eq:Q85:7.8}
\]
The bounding quadratic decreases on \([1/4,\allowbreak 1/3]\).
Thus \(T(f_2h_4)\le\allowbreak 205/256\) and
\(pT^2(f_2h_4)\le\allowbreak 205/256\).
Its eighth power is smaller than \(e^{-d/3}\):
\((205/256)^2<\allowbreak 13/20\) gives an eighth power below \(9/50\), whereas
\(e^{-d/3}\ge\allowbreak 8/9\).
For the canonical pair, its exact mean gives
\[
 E(h_2f_4)=1-3d+\delta_2-d^2+\delta_3
 \le1-(2487/1024)d-d^2\le1353/4096<1/3.              \tag{Q85.7.9}\label{eq:Q85:7.9}
\]
Hence \(pT^3(h_2f_4)\le\allowbreak (2/3)p\bar f_2\).
Equation \eqref{eq:Q85:7.6} gives an eighth moment at most
\((2/3)^8(82/15)<\allowbreak 1/2<\allowbreak e^{-d/3}\).
Together with \eqref{eq:Q85:7.2} and \eqref{eq:Q85:7.5}, this proves
\eqref{eq:Q85:0.2} in its fixed directions.

\paragraph{Q85.8 Every \(b\ge\allowbreak 5,\allowbreak L=\allowbreak 3\) on the low domain}
Let \(0<\allowbreak d\le\allowbreak 1/5,\allowbreak F=\allowbreak h_2f_b\). Three-step Duhamel propagation of the actual square gives
\[
 p+Z=2+\varepsilon-d-k\rho+N,
\]
\[
 N=K'V_3,\quad
 V_3=\varepsilon1+(1+\delta)g'+K'g'+kK'^2R,
\]
\[
 EN=-(1+\delta)d+kD_3,\qquad D_3=\int x\mathsf A.
\]
Consequently,
\[
 \log EH^8\le16\log m-16d+8\varepsilon-8\delta d
                 +8k(D_3-\rho)+17\|V_3\|^2.          \tag{Q85.8.1}\label{eq:Q85:8.1}
\]
The nonconstant coordinates of \(V_3\) are
\(1+\delta-kx\mathsf A\) and \(1+kx\mathsf B\).
For every \(b\ge\allowbreak 5\), their squared norm satisfies
\[
 (1+\delta-kx\mathsf A)^2+x(1+kx\mathsf B)^2
 \le(1+\delta)^2+x.                                 \tag{Q85.8.2}\label{eq:Q85:8.2}
\]
Here is the full finite certificate.
Using \(\mathsf A^2+x\mathsf B^2\le\allowbreak  u_0\mathsf A+xz_0\mathsf B\) and
\(\mathsf A\ge\allowbreak (1-x)\mathsf B\), the unshifted squared difference is at most
\[
 kx\mathsf B[-2+4x+kx\{u_0(1-x)+xz_0\}].
\]
The substitution is valid because \(2-kxu_0>\allowbreak 0\).
The bracket is negative on the full rectangle:
\(x[u_0(1-x)+xz_0]\) increases with \(x\), with derivative at least
\(u_0(1-2x)>\allowbreak 0\), and all parameters increase with \(d\).
At \(d=\allowbreak 1/5,\allowbreak x=\allowbreak 13/32\), its negative margin is \(79745/3174400\).
A nonnegative \(\delta\) shift adds \(-2\delta kx\mathsf A\) to the difference.
This proves \eqref{eq:Q85:8.2}, including \(b=\allowbreak 5\).
The constant coordinate then gives
\[
 \|V_3\|^2\le(1+\delta)^2d+\delta_2
                 +\varepsilon^2-2\varepsilon(d+k\rho).          \tag{Q85.8.3}\label{eq:Q85:8.3}
\]

Besides \(\varepsilon\le\allowbreak 3d/5\), the parameter and endpoint bounds give
\[
 -\varepsilon\le k(a_1-D_0)<7d/20,\qquad k\rho<3d/10.
\]
The full \(\varepsilon\) contribution satisfies
\[
 8\varepsilon+17\varepsilon^2-34\varepsilon(d+k\rho)
 \le(24/5)d-(357/25)d^2.                             \tag{Q85.8.4}\label{eq:Q85:8.4}
\]
For \(\varepsilon\ge\allowbreak 0\), discard the nonpositive \(\rho\) term.
The convex quadratic on \([0,\allowbreak 3d/5]\) has maximum at an endpoint;
\(8-34d>\allowbreak 0\) makes the upper endpoint larger, giving exactly the right side.
For \(\varepsilon\le\allowbreak 0\), put \(w=\allowbreak -\varepsilon\le\allowbreak 7d/20\).
The upper bound is a convex quadratic in \(w\), with endpoints zero and
\[
 -(14/5)d+(7021/400)d^2.
\]
The right side of \eqref{eq:Q85:8.4} is positive and exceeds this endpoint:
their difference is \(d[38/5-(12733/400)d]>\allowbreak 0\) for \(d\le\allowbreak 1/5\).
Thus the actual mean of either sign and its constant-vector cross term remain included.
Use \(\delta\le\allowbreak 6d/5\), \eqref{eq:Q85:4.9}, and
\[
 16\log m\le-18d-(81/8)d^2,
 \quad 8k\theta^4<3/10,\quad8k\theta<17/4.
\]
Substitution into \eqref{eq:Q85:8.1}--\eqref{eq:Q85:8.4} gives
\[
 \log EH^8\le-(799/160)d+(2209/200)d^2+(612/25)d^3
 \le-(36111/20000)d<-(9/5)d.                         \tag{Q85.8.5}\label{eq:Q85:8.5}
\]
All remaining positive coefficients were bounded at \(d=\allowbreak 1/5\).
This proves the canonical low-domain row for every finite \(b\ge\allowbreak 5\).

\paragraph{Q85.10 Completion using the actual long-connector inputs}
For the canonical \(F=\allowbreak h_2f_b\), Section~\ref{proof:Q5-C086} proves
\(\E[pT^LF]^8\le\allowbreak 1-d/20000\) when \(b=\allowbreak 3,\allowbreak L\ge\allowbreak 3\), or \(b\ge\allowbreak 4,\allowbreak L\ge\allowbreak 4\).
It also proves, for every finite \(b\ge\allowbreak 5\) and \(1/5\le\allowbreak  d\le\allowbreak 1/3\),
\(\E(h_2f_b)\le\allowbreak (1-d)/2\).

For the remaining \(b\ge\allowbreak 5,\allowbreak L=\allowbreak 3,\allowbreak d\ge\allowbreak 1/5\) row, the kernel
\(0\le\allowbreak 1-K(x,\allowbreak y)\le\allowbreak 2\) gives \(TF\le\allowbreak 1-d\). Positivity yields
\[
 pT^3F\le(1-d)p\bar f_2.
\]
The actual path estimates \eqref{eq:Q85:7.6}--\eqref{eq:Q85:7.7} now give
\[
 E[pT^3F]^8\le
 \begin{cases}
 (4/5)^8(35/8)=286720/390625<3/4,&1/5\le d\le1/4,\\
 (3/4)^8(82/15)=538002/983040<3/4,&1/4\le d\le1/3.
 \end{cases}                                                   \tag{Q85.10.1}\label{eq:Q85:10.1}
\]
Since \(e^{-d/5}\ge\allowbreak 1-d/5\ge\allowbreak 14/15\), this proves the second inequality in
\eqref{eq:Q85:0.3} on the remaining domain.
Together with \eqref{eq:Q85:8.5}, it proves that row in the same direction.

The full ordinary catalogue has the following rows:
\[
\begin{array}{c|c|c}
\text{integer row}&F_{b,L}&\text{ordinary eighth upper bound}\\ \hline
b\ge3,\ L=1&h_2f_b&1-d/500\\
b=3,\ L=2&h_2f_3&e^{-d/25}\\
b=3,\ L\ge3&h_2f_3&1-d/20000\\
b=4,\ L=2&f_2h_4&e^{-d/3}\\
b=4,\ L=3&h_2f_4&e^{-d/3}\\
b\ge5,\ L=2&f_2h_b&e^{-d/5}\\
b\ge5,\ L=3&h_2f_b&e^{-d/5}\\
b\ge4,\ L\ge4&h_2f_b&1-d/20000
\end{array}
\]
The unit row uses the \((2,\allowbreak 3)\) unit input of Section~\ref{proof:Q5-C053-unit} and
Section~\ref{proof:Q5-C057} for \(b\ge\allowbreak 4\).
The \(b=\allowbreak 3,\allowbreak L\ge\allowbreak 3\) and \(b\ge\allowbreak 4,\allowbreak L\ge\allowbreak 4\) rows use Section~\ref{proof:Q5-C086};
all other rows were proved above.
For \(0\le\allowbreak  d\le\allowbreak 1/3\),
\[
 e^{-d/25}\le\frac1{1+d/25}
 =1-\frac{d}{25+d}\le1-\frac{3d}{76}
 \le1-d/20000.
\]
The stronger short-row exponentials satisfy the same bound. Therefore
\[
 \boxed{E[pT^L F_{b,L}]^8\le Q(d)=1-d/20000
   \quad(b\ge3,L\ge1).}                              \tag{Q85.10.2}\label{eq:Q85:10.2}
\]

\paragraph{Q85.11 Decorations and complete reversal}
At the terminal, let \(E_t\) be the product of its original equal pairs.
One member of each pair points entirely outward; the other has its first edge inward and its remaining edges outward.
At the receiving root, let \(E_r\) be the corresponding product of actual pairs or their complete reversals, as specified by the structure.
Each pair factor lies in \([0,\allowbreak 1]\).
The decorated actual message is
\[
 H_D=E_rpT^L(E_tF_{b,L}).                              \tag{Q85.11.1}\label{eq:Q85:11.1}
\]
Positivity and \eqref{eq:Q85:10.2} imply \(\E H_D^8\le\allowbreak  Q(d)\);
the decorations are inserted before their outgoing propagation.

For complete reversal, replace \(K\) by \(-K\) on the entire fixed structure.
Then \(p,\allowbreak q\) exchange, \(A,\allowbreak T\) exchange, \(f_n,\allowbreak \bar f_n\) exchange, and
\(h_n\) becomes \(A\bar f_{n-1}\).
All terminal and root decorations reverse together.
In particular, the \(L=\allowbreak 2,\allowbreak b\ge\allowbreak 4\) row reverses \(f_2h_b\) as a whole.
The reversed ordinary bound has the same \(Q(d)\) and the same original edge and vertex counts.
